\documentclass[a4paper]{article}
% \usepackage[francais]{babel}
\usepackage[utf8]{inputenc}
\usepackage[T1]{fontenc}
\usepackage{datetime}
\usepackage{lmodern}
\usepackage{hyphenat}
\usepackage{microtype}
\usepackage{newspaper}
\usepackage{hyperref}
\usepackage{graphicx}
\usepackage{multicol}
\usepackage{picinpar}
\usepackage{newspaper-mod}
\usepackage{lipsum}
\usepackage{epigraph}
\usepackage{blindtext}
\usepackage[footnote,printonlyused]{acronym}
\usepackage{scrextend}
\usepackage{bookmark}
\usepackage[textwidth=20cm,textheight=28.7cm]{geometry}

\usepackage{graphicx}
\graphicspath{ {assets/images/} }

\usepackage[
  backend=bibtex,
  sorting=none
]{biblatex}
\addbibresource{bibliographie.bib}
\AtNextBibliography{\footnotesize}

\setlength{\footnotesep}{0cm}
\setlength{\skip\footins}{0cm}
\date{Le vendredi 03 septembre 2021}
\currentvolume{2}
\currentissue{2}
\SetPaperName{Le Glas:en}
\SetHeaderName{Le Glasen}
\SetPaperLocation{France Libre}
\SetPaperSlogan{Sonnons le glas de la soumission et de l'indignité}
\SetPaperPrice{Hebdomadaire de Réinformation Gratuit}
\thispagestyle{empty}
\pagenumbering{gobble}

\deffootnote[1.0em]{3.2em}{10em}
{\textsuperscript{\thefootnotemark}\,\enskip}

\begin{document}

\begin{multicols}{3}

  \begin{center}
    \vspace*{1cm}
            {\textgoth{\fontsize{15}{60}\selectfont\usefont{LYG}{bigygoth}{m}{n}Le Glas:en}}\hfill \\
            \vspace{0.1cm}
            \footnotesize{Sonnons le Glas de la Soumission}
  \end{center}
  
  \vspace*{0.1in}
  \begin{center}
    \MakeUppercase{\small 03 septembre 2021 \\ Volume n°2}\hfill \\
    \vspace*{0.2in}

    \begin{minipage}{5cm}
      \begin{center}
        \noindent «La philosophie nous enseigne à douter de ce qui nous
        paraît évident. La propagande, au contraire, nous enseigne à
        accepter pour évident ce dont il serait raisonnable de douter.»
        -- Aldous Huxley
      \end{center}
    \vspace*{1cm}
    \end{minipage}
    
    % \vspace*{0.1in}
    \vspace*{-16pt}
    \rule[1pt]{\hsize}{1pt}
  \end{center}

  \headline{Traitement Génique}

  La pharmacovigilance de l'Union Européenne ADR Reports
  \cite{AdrReports}, regroupant les remontées des pays membres, met à
  jour ses données tous les samedi matin. En date du samedi 28 août
  2021, \textbf{23252 décès} ont été déclarés ainsi que \textbf{580050
    effets indésirables non résolus}, \textbf{904943 effets
    indésirables résolus}, \textbf{26667 effets indésirables résolus
    avec séquelles}, \textbf{501802 effets indésirables en cours de
    résolution} et \textbf{373916 effets indésirables dont le statut
    est inconnu}. Au moment de l'extraction de ces données, selon le
  site Vaccine Tracker \cite{VaccineTracker}, environ 281 millions
  d'européens ont reçu une première injection, et environ 250 millions
  d'européens ont reçu deux injections.

  Les rapports d'effets indésirables aux États-Unis se font au moyen
  du système VAERS \cite{OpenVaers}, dont la base de données est
  alimentée par la FDA et le CDC. En date du 20 août 2021,
  \textbf{623341 déclarations} ont été remontées, dont \textbf{13627
    décès}, \textbf{55821 hospitalisations}, \textbf{74368
    hospitalisations en urgences}, \textbf{100966 visites médicales},
  \textbf{5721 chocs anaphylactiques}, \textbf{4721 paralysies de
    Bell}, \textbf{1671 fausse-couches}, \textbf{6071 arrêts
    cardiaques}, \textbf{5093 myocardites}, \textbf{17794 invalides à
    vie}, \textbf{2831 cas de thromboses}, \textbf{25921 réactions
    allergiques graves} et \textbf{7324 zonas}. Au 31 août, environ
  205 millions de personnes ont reçu une dose et environ 174 millions
  ont reçus 2 injections \cite{UsaFacts}.

  Au Royaume-Uni, les données sont remontées au moyen du système
  Yellow Card\cite{YellowCardReport}. Les derniers rapports datent du
  18 août 2021, et n'ont pas évolué depuis le dernier numéro. Pour
  rappel, environ 46 millions de premières doses ont été injectées et
  environ 40 millions de citoyens britanniques ont reçu une seconde
  injection. \textbf{302146 effets indésirables} dont \textbf{506
    décès} ont été remontés pour le produit de Pfizer, \textbf{229134
    effets indésirables} dont \textbf{1056 décès} ont été remontés
  pour le produit issu d'Astrazeneca, \textbf{14019 effets
    indésirables} dont \textbf{17 décès} ont été déclarés pour le
  produit de chez Moderna, et \textbf{1036 remontés} dont \textbf{28
    décès} pour des produits d'origine inconnue.

  L'Australie reporte les effets indésirables au moyen du DAEN
  \cite{Australia:DAEN} pour 3 traitements géniques (Pfizer,
  Astrazeneca, Moderna). En date du 17 août 2021, \textbf{49570
    déclarations} ont été remontées dont \textbf{465 décès}. Dans ces
  rapports, outre les réactions “classiques”, 1163 douleurs
  thoraciques, 1632 cas de douleurs aux extrémités, 716 cas d'embolies
  pulmonaires, 778 cas thromboses veineuses, 633 cas de vision
  troublées ou encore 445 chocs anaphylactiques.

  Le Medeffect \cite{Canada:CVAROD} est en charge de la collecte des
  effets indésirables au Canada. En date du 28 aout 2021, cette entité
  a reçu \textbf{649 déclarations}, dont \textbf{15 décès} pour
  environ 53 millions de doses administrées. La pharmacovigilance de
  l'Afrique du Sud se nomme quant à elle AEFI \cite{SouthAfrica:AEFI}
  et est gérée par la SAHPRA.  En date du 30 août 2021, environ 12
  millions de doses ont été administrées, soit 2,7 millions de doses
  du vaccin Janssen et 9,5 millions de doses du vaccin Pfizer. Les
  derniers rapports de l'AEFI datent du 13 août 2021, \textbf{1932
    effets indésirables} ont été déclarés dont \textbf{58 décès
    suspectés}.

  À noter que les chiffres présentés ici sont à mettre en
  perspective. Entre \textbf{1 et 10\% des déclarations sont
    réellement remontées} \cite{Study:UnderReporting} dans les données
  de pharmacovigilance par les différents acteurs. De plus, ces
  traitements, après avoir passé les phases 1 et 2, sont désormais en
  phase 3 des tests. Cette phase correspond à un test grandeur nature
  se déroulant sur plusieurs années. Ce test est important pour savoir
  si le produit est fonctionnel, a un effet positif sur l'immunité,
  mais aussi pour s'assurer qu'il n'y ait pas d'effets indésirables
  sur le long terme. Plusieurs questions restent en suspens, tout
  d'abord cette phase de test doit généralement inclure des placebos
  pour qu'il y ai un groupe témoin, puis, selon les protocoles
  utilisés, différents dosages peuvent aussi être utilisés en fonction
  des populations ciblées.

  \closearticle

  \headline{Entre Corruption et Tyrannie Sanitaire}

  Le peuple français et plus largement les citoyens vivant en Occident
  seraient-ils devenu les victimes d'une expérience sociale de grande
  ampleur? Entre expérience de Milgram et de Stanford, une chose est
  certaine, la population est de plus en plus divisée par la caste
  politico-médiatique dont les dérives douteuses s'invitent de force
  dans la vie de tout à chacun. D'autres entités sont aussi à blâmer,
  telles que les multinationales pharmaceutiques, plus communément
  surnommées “Big Pharma”. Ces nouveaux profiteurs de guerres ne sont
  plus à une action crapuleuse près. Prises de nombreuses fois la main
  dans le sac, et ce malgré des sanctions records, la prise de risque
  reste bien trop souvent en faveur de l'exaction. Effectivement, le
  profit étant largement supérieur à la punition, un simple calcul
  entre le bénéfice et le risque encouru ne laisse aucun doute sur le
  résultat final. Voici une courte liste non-exhaustive décrivant le
  comportement de ces structures.

  La société Moderna fut créée en 2010 dans le Massachussets aux
  États-Unis. Elle est spécialisée dans la conception de vaccins
  utilisant la technologie ARN messager. Créée récemment, elle n'a
  aujourd'hui aucune infraction à son actif, mais cela ne saurait
  tarder. \textbf{Le japon a dû supprimer 2,6 millions de doses}
  \cite{NewsWeek:JapanModerna} du traitement génique Moderna après
  avoir trouvé des particules métalliques issue d'une probable
  contamination dans un lot.

  La société Astrazeneca, fondée en 1999, est le résultat la fusion de
  la société suédoise Astra AB créée en 1913 par Knut Sjöberg
  (pharmacien) et la société anglaise Zeneca conçue en 1993. Elle se
  spécialise dans les produits pharmaceutiques et la
  biotechnologie. En 2010, Astrazeneca dû payer \textbf{520 millions
    de dollars} pour avoir testé un neuroleptique
  \cite{Justice:Astrazeneca2010} sur des personnes âgées, des vétérans
  et des prisonniers, devenus alors cobayes. En 2016, elle dû payer
  \textbf{5,5 millions de dollars} pour avoir corrompu des médecins en
  Chine et en Russie. Le site GoodJobFirst
  \cite{GoodJobFirst:Astrazeneca} recense \textbf{21 infractions} pour
  environ \textbf{1,1 milliards de dollars de pénalité} (comptabilisé
  depuis 2000).

  La société Janssen fut créée en 1953 en Belgique par Paul Janssen
  (chimiste/pharmacologue). Rachetée par le conglomérat Johnson \&
  Johnson en 1961, elle rentre alors dans la branche Recherche et
  Développement de cette structure. Elle se spécialise sur un large
  spectre de trouble de la santé humaine. En 2013, elle dû payer
  \textbf{2,2 milliards de dollars} pour pot-de-vins et médicaments
  sans étiquettes \cite{Justice:JohnsonJohnson2013}. En 2019, elle
  doit payer \textbf{465 millions de dollars} pour avoir contribué à
  la crise des Opioïdes\cite{Justice:JohnsonJohnson2019}, puis
  \textbf{116 millions de dollars} pour publicités
  mensongères\cite{Justice:JohnsonJohnson2019-2}. Le site GoodJobFirst
  \cite{GoodJobFirst:JohnsonJohnson} recense \textbf{58 infractions},
  pour environ \textbf{4,2 milliards de dollars} de pénalités
  (comptabilisé depuis 2000).

  La société Pfizer, créée en 1849 à New York par les businessmen
  Charles Pfizer et Charles F. Erhart, est spécialisée dans la
  création de médicaments et de vaccins. Elle a eu le privilège de
  recevoir l'accord en justice le plus élevé de l'histoire des
  États-Unis en 2009: \textbf{2,3 milliards de dollars} pour marketing
  frauduleux \cite{Justice:Pfizer2009}. En 2016, elle a dû payer
  \textbf{784 millions de dollars} pour des affaires de fraudes sur
  les prix\cite{Justice:Pfizer2016}. En 2013, elle a dû payer
  \textbf{491 millions de dollars} pour de la commercialisation
  illégale de médicaments pour des utilisations non
  approuvées\cite{Justice:Pfizer2013}. En 2006, elle dû payer des
  compensations à des familles nigériennes durant la pire épidémie de
  méningite africaine, causant la mort de 100 enfants et des dommages
  au foie\cite{Justice:Pfizer2006}. Le site GoodJobFirst
  \cite{GoodJobFirst:Pfizer} recense \textbf{71 infractions}, pour
  environ \textbf{4,6 milliards de dollars de pénalités} (comptabilisé
  depuis 2000).

  \closearticle

  \headline{Défense Juridique}

  L'actualité juridique a été très chargée. Maitre Carlo Alberto
  Brusa, président de l'association Reaction19, a annoncé lors d'un
  live \cite{Reaction19:JuridiqueNommage} le lancement d'une nouvelle
  office juridique. Cette dernière est axée sur un problème de
  définition des produits en cours d'utilisation. Effectivement, le
  nom de ces produits est sensé être standard et partagé sur tous les
  documents officiels, suivant une nomenclature stricte, ce qui n'est
  malheureusement pas le cas aujourd'hui.

  Le 27 août 2021, le collectif Reinfo-Liberté a annoncé porter
  plainte contre X \cite{ReinfoLiberte:PlainteX} pour faux et usage de
  faux, escroquerie et trafic d’influence. Cette plainte vise à
  dénoncer un article scientifique de l'Institut Pasteur repris par le
  Conseil Scientifique pour justifier la mise en place du ``passe
  sanitaire''. Le même jour est annoncé que le pass sanitaire n'est
  plus obligatoire dans les centres commerciaux d'Ile de France
  \cite{CourrierDuSoir:SuspensionPassIdf}, suite à un ordre de
  suspension effectué par les avocats Yoanne Sibille et Tarek
  Koraitem. Quelques jours après, le même ordre est arrivé jusqu'en
  Haut-Rhin \cite{SudOuest:SuspensionPassHautRhin}. Le gouvernement a
  réagi en faisant appel \cite{}. Affaire à suivre de très près.

  Le 31 août 2021 se tenait une audience des plus importantes au
  tribunal de grande instance de Paris. Maitre Diane Protat plaidait
  contre le pass sanitaire et l'obligation vaccinale
  \cite{FranceSoir:MeProtatAudience}. Trois magistrats ont été nommés,
  fait extrêmement rare pour ce genre d'affaires, qui montre l'ampleur
  de la situation. Maitre Protat résume son action et son plan
  d'attaque dans un article paru dans le journal France Soir
  \cite{FranceSoir:MeProtatDebrief}.

  Le 1er septembre 2021 voyait aussi la naissance du Syndicat Liberté
  Santé (SLS) \cite{SyndicatLiberteSante:Naissance}. Ce nouveau
  syndicat a été créé pour défendre les droits des professionnels de
  santé en suivant une ligne de conduite stricte tout en contribuant
  ”à la construction d’une société plus équitable, riche de ses
  individualités, dans la paix, la prudence, la communication loyale,
  l’écoute attentive, la discussion, et en ne nous laissant pas
  acheter à des fins opposées à nos valeurs, mais restant ouvert à des
  contributions bienveillantes''.

  Le collectif ReinfoCovid met aussi à disposition une boite à outils
  \cite{ReinfoCovid:BoiteOutil} permettant de vous donner différentes
  options pour votre défense. Par ailleurs, soyez toujours prudent,
  soyez toujours alerte et posez vous des questions
  \cite{FloraHebridge:Conseil}. Ne prenez jamais pour acquis tout ce
  qui est partagé, vérifiez les sources et cherchez par vous même. La
  lutte ne fait que commencer. Courage à vous tous.
  
  \closearticle

  \headline{Témoignages}

  De nombreux témoignages en provenance de toute la France sont
  partagés sur le réseau internet, par le bouche à oreille ou tout
  simplement lors des manifestations et de nos nombreuses
  rencontres.

  Un pays ne se résume pas qu'à des chiffres sans âmes. Toutes ces
  expériences que vous vivez mettent en avant les souffrances et les
  inquiétudes vis à vis de la situation actuelle. Pour vous permettre
  de vous faire entendre, nous vous proposons l'ouverture des colonnes
  du Glasen et ainsi pouvoir partager votre vécu.
  
  \closearticle
  
\begin{center}
  \rule[1pt]{4.5cm}{1pt}
  \textbf{\url{https://t.me/leglasen}}
  \includegraphics[width=4.5cm,height=4.5cm]{qrcode-telegram.eps}
  \textbf{\href{mailto:leglasen@protonmail.com}{leglasen@protonmail.com}} \\
  \vspace*{0.2cm}
  \footnotesize\textbf{\href{https://creativecommons.org/licenses/by-nc-nd/4.0/legalcode.fr}{CC BY-NC-ND 4.0}}
  \vspace*{1cm}
  \rule[1pt]{4.5cm}{1pt}
\end{center}

\printbibliography[heading=none]

\end{multicols}
\end{document}
